\documentclass[10pt,a4paper]{amsart}
\usepackage[margin=19mm]{geometry}
\usepackage{amsmath,amssymb,mathtools}
\usepackage[scr=rsfs,cal=euler]{mathalfa}
\usepackage{xfrac}
\usepackage[unicode,hidelinks,bookmarks=false]{hyperref}
\newcommand{\ie}{i.e.\ }
\newcommand{\cf}{cf.\ }
\newcommand{\resp}{respectively\ }
\newcommand{\Subsection}{\subsection}
\providecommand{\nobreakdash}{\nobreak\hbox{-}}
\setlength{\emergencystretch}{2em}
\multlinegap0pt
\usepackage{tikz}
\usepackage{csquotes}
\usepackage{float}
\usepackage{xcolor}
\usepackage[shortlabels]{enumitem}
\usepackage{amssymb}
\newtheorem{proposition}{Proposition}
\newtheorem{theorem}{Theorem}
\title{A Modular Form Proof of the Irrationality of $\zeta\left(3\right)$}
\author{Pang Ern Thang}
\date{Source: arXiv:2607.17123v1 (CC BY 4.0)}
\begin{document}
\maketitle

\noindent\emph{Source and licence.} A Modular Form Proof of the Irrationality of $\zeta\left(3\right)$. Version arXiv:2607.17123v1 (CC BY 4.0), distributed by arXiv under the Creative Commons Attribution 4.0 International licence, http://creativecommons.org/licenses/by/4.0/, as recorded in the arXiv licence metadata for this exact version and shown on the abstract page https://arxiv.org/abs/2607.17123v1. Full licence notice: \texttt{licenses/arxiv-2607.17123-CC-BY-4.0.txt}.\par\smallskip

\noindent\textit{Editorial scope.} Counted unit: the article's theorem theorem (source0.tex lines 712--714). The article's complete original proof of it is delivered with it (source0.tex lines 715--750, one \texttt{proof} environment of 36 lines). No mathematics is written, repaired or re-derived: the statement and the proof are the article's own lines, in the article's own notation, and no retained line is substituted or re-flowed.  Retained uncounted as the source-local context the proof needs: lines 113--115; lines 127--129; lines 393--395.  Nothing was omitted from any retained span, so the package carries no omission record.  No retained line contains a citation, so the wrapper carries no bibliography and no bibliography entry.  No retained line contains a figure environment, an image-inclusion command or drawing code, so no image is delivered and the package declares no asset dependency.  Editorial formatting only, in addition: an AMS \texttt{amsart} preamble that restates the article's own declarations and macros used by the retained lines, namely the environments \texttt{proposition}, \texttt{theorem} on separate counters, together with the standard packages those lines need.  The retained environments print in this document as Proposition 1, Theorem 1 in order of appearance, whereas the article prints the same items with the numbering of its own full document.  The complete native source of the article as distributed in the arXiv e-print for this exact version is bundled unchanged as \texttt{sources/205537/source0.tex}.
\begin{proposition}[Beukers' irrationality criterion]\label{prop: Beukers' irrationality criterion}
Let $f_0\left(t\right),f_1\left(t\right),\ldots,f_k\left(t\right)$ be power series in $t$. Suppose for any $n\in\mathbb{N}$ and $i=0,1,\ldots,k$, the $n^\text{th}$ coefficient in the Taylor series of $f_i$ is rational and has denominator dividing $d^n\left(\operatorname{lcm}\left(1,\ldots,n\right)\right)^r$, where $r,d\in\mathbb{N}$ are fixed. Suppose there exist $\theta_1,\ldots,\theta_k\in\mathbb{R}$ such that $f_0\left(t\right)+\theta_1f_1\left(t\right)+\cdots+\theta_kf_k\left(t\right)$ has radius of convergence $\rho$ and infinitely many non-zero Taylor coefficients. If $\rho>de^r$, then at least one of $\theta_1,\ldots,\theta_k$ is irrational.
\end{proposition}

\begin{proposition}\label{prop: h formula in terms of d}
Let $q=e^{2\pi i\tau}$ and \[F\left(\tau\right)=\sum_{n=1}^{\infty}a_nq^n\] be a Fourier series convergent for $\left|q\right|<1$ such that for some $k,N\in\mathbb{N}$ and $\varepsilon\in\left\{-1,1\right\}$, \[F\left(-\frac{1}{N\tau}\right)=\varepsilon\left(-i\tau\sqrt{N}\right)^kF\left(\tau\right).\] Let \[f\left(\tau\right)=\sum_{n=1}^{\infty}\frac{a_n}{n^{k-1}}q^n.\] Let \[L\left(F,s\right)=\sum_{n=1}^{\infty}\frac{a_n}{n^s}\] be the associated Dirichlet series and define \[h\left(\tau\right)=f\left(\tau\right)-\sum_{0\le r<\frac{k-2}{2}}\frac{L\left(F,k-r-1\right)}{r!}\left(2\pi i\tau\right)^r.\] Then, \[h\left(\tau\right)-D=\left(-1\right)^{k-1}\varepsilon\left(-i\tau\sqrt{N}\right)^{k-2}h\left(-\frac{1}{N\tau}\right),\] where \[D=\begin{cases}0&\text{if $k$ is odd};\\[4pt]\displaystyle\frac{L\left(F,\frac{k}{2}\right)}{\left(\frac{k}{2}-1\right)!}\left(2\pi i\tau\right)^{\frac{k}{2}-1}&\text{if $k$ is even}.\end{cases}\]
\end{proposition}

denote the modular discriminant, which is a cusp form of weight 12 for $\operatorname{SL}_2\left(\mathbb{Z}\right)$. Then, \begin{align}\label{eqn: beukers formula for t}
t\left(\tau\right)=\left(\frac{\Delta\left(6\tau\right)\Delta\left(\tau\right)}{\Delta\left(3\tau\right)\Delta\left(2\tau\right)}\right)^{1/2}=\left(\frac{\eta\left(6\tau\right)\eta\left(\tau\right)}{\eta\left(3\tau\right)\eta\left(2\tau\right)}\right)^{12}=q\prod_{n=0}^{\infty}\left(1-q^{6n+1}\right)^{12}\left(1-q^{6n+5}\right)^{12}.
\end{align}

\begin{theorem}[Apéry]
$\zeta\left(3\right)$ is irrational.
\end{theorem}

\begin{proof}
Define the functions $F$ and $E$ satisfying \begin{align*}
    40F\left(\tau\right)&=E_4\left(\tau\right)-36E_4\left(6\tau\right)-28E_4\left(2\tau\right)+63E_4\left(3\tau\right)\\
    24E\left(\tau\right)&=-5E_2\left(\tau\right)+30E_2\left(6\tau\right)+2E_2\left(2\tau\right)-3E_2\left(3\tau\right)
\end{align*}
The function \(F\) belongs to \(S_4\left(\Gamma_1\left(6\right)\right)\), the complex vector space of cusp forms of weight 4 for the congruence subgroup $\Gamma_1\left(6\right)$, and satisfies the Fricke transformation law
\[F\left(-\frac{1}{6\tau}\right)=-36\tau^4F\left(\tau\right).\]
The Dirichlet series corresponding to $F\left(\tau\right)$ is \begin{align*}
L\left(F,s\right)&=\sum_{n=1}^{\infty}\left[\frac{6\sigma_3\left(n\right)}{n^s}-36\frac{6\sigma_3\left(n\right)}{\left(6n\right)^s}-28\frac{6\sigma_3\left(n\right)}{\left(2n\right)^s}+63\frac{6\sigma_3\left(n\right)}{\left(3n\right)^s}\right]\\
&=6\left(1-6^{2-s}-7\cdot 2^{2-s}+7\cdot 3^{2-s}\right)\zeta\left(s\right)\zeta\left(s-3\right)
\end{align*}
where we used the fact that \[\sum_{n=1}^{\infty}\frac{\sigma_k\left(n\right)}{n^s}=\zeta\left(s\right)\zeta\left(s-k\right)\quad\text{where }\operatorname{Re}\left(s\right)>k+1.\]
One can use the functional equation for the Riemann zeta function \[\zeta\left(s\right)=2^s\pi^{s-1}\sin \left(\frac{\pi s}{2}\right)\Gamma\left(1-s\right)\zeta\left(1-s\right)\]
and the gamma function $\Gamma\left(s+1\right)=s\Gamma\left(s\right)$ to show that $\zeta\left(0\right)=-\frac{1}{2}$. Consequently, \[L\left(F,3\right)=\zeta\left(3\right).\]
Let \(f\) be the Fourier series normalised by
\[f'''\left(\tau\right)=\left(2\pi i\right)^3F\left(\tau\right)\quad\text{where }
f\left(i\infty\right)=0.
\]
It follows from Proposition \ref{prop: h formula in terms of d} that
\begin{align}\label{eqn: transformation of f}
6\tau^2\left(f\left(-\frac{1}{6\tau}\right)-L\left(F,3\right)\right)=-f\left(\tau\right)+L\left(F,3\right)-2\pi i\tau L\left(F,2\right).
\end{align}
One can show that $L\left(F,2\right)=0$, so \[6\tau^2\left(f\left(-\frac{1}{6\tau}\right)-\zeta\left(3\right)\right)=\zeta\left(3\right)-f\left(\tau\right).\]
Set \[H\left(\tau\right)=E\left(\tau\right)\left(f\left(\tau\right)-\zeta\left(3\right)\right).\]
One can show that $E$ satisfies the Fricke transformation law \[E\left(-\frac{1}{6\tau}\right)=-6\tau^2E\left(\tau\right)\quad\text{so}\quad H\left(-\frac{1}{6\tau}\right)=H\left(\tau\right).\]
Recall our formula for $t$ in (\ref{eqn: beukers formula for t}) obtained from the product expansion of $\eta$. Then, $t=q+O\left(q^2\right)$ so $t$ has a local compositional inverse $q=q\left(t\right)\in t\mathbb{Z}\left[\left[t\right]\right]$ near $t=0$. As such, we regard $E,f,H$ as power series in $t$. 

We previously established that the first two positive branching values are $t_1=(\sqrt{2}-1)^4$ and $t_2=(\sqrt{2}+1)^4$. The first branching value corresponds to the fixed point $\tau=\frac{i}{\sqrt{6}}$ of the Fricke involution. The two local branches of the inverse $\tau=\tau\left(t\right)$ near $t=t_1$ are interchanged by $\tau \mapsto -\frac{1}{6\tau}$. Note that the function $H$ takes the same value on these two branches. Hence, the apparent branch singularity at $t=t_1$ is removable for $H$. The next branching value is $t_2$, and therefore the Taylor series of $H\left(t\right)$ about $t=0$ has radius of convergence $\rho=t_2=17+12\sqrt{2}$. In particular, this radius is finite, so $H\left(t\right)$ is not a polynomial and has infinitely many non-zero Taylor coefficients.

It remains to control the denominators. Note that the Fourier coefficients $a_n$ of $F$ are integers. So, consider the Eichler integral of $F$ by \[f\left(\tau\right)=\sum_{n=1}^{\infty}\frac{a_n}{n^3}q^n\]
which has rational coefficients whose $n^\text{th}$ denominator divides $n^3$. Let $D_n=\operatorname{lcm}\left(1,2,\ldots,n\right)$. SInce $q\left(t\right)\in t\mathbb{Z}\left[\left[t\right]\right]$, the coefficient of $t^n$ in $f\left(q\left(t\right)\right)$ is an integral linear combination of $\frac{a_m}{m^3}$, where $1\le m\le n$. Since $m\mid D_n$ for all $m\le n$, then its denominator divides $D_n^3$. 

Likewise, $E\left(q\right)\in \mathbb{Z}\left[\left[q\right]\right]$ so $E\left(t\right)\in\mathbb{Z}\left[\left[t\right]\right]$. Write \[E\left(t\right)f\left(t\right)=\sum_{n=0}^{\infty}A_nt^n\quad\text{and}\quad E\left(t\right)=\sum_{n=0}^{\infty}B_nt^n.\]
So, $A_n\in \frac{1}{D_n^3}\mathbb{Z}$ and $B_n\in\mathbb{Z}$ and \[H\left(t\right)=\sum_{n=0}^{\infty}\left(A_n-\zeta\left(3\right)B_n\right)t^n.\]
Apply Proposition \ref{prop: Beukers' irrationality criterion} with $f_0\left(t\right)=E\left(t\right)f\left(t\right)$, $f_1\left(t\right)=-E\left(t\right)$, and $\theta=\zeta\left(3\right)$. Also, set $d=1$ and $r=3$. Then, $H\left(t\right)$ has infinitely many non-zero coefficients and $\rho=17+12\sqrt{2}>e^3$. This implies that $\zeta\left(3\right)$ is irrational.
\end{proof}

\end{document}
